\section{Finite evidence and reproduction}\label{sec:certificates}
The accompanying computations distinguish three assertions: a literal
record satisfies its algebraic equations; a stated finite domain is
covered by such records; and an infinite counting theorem follows from
that finite domain together with proved structural and summation
lemmas. These assertions require different arguments.

The publication package contains the following mathematical data.
\begin{center}
\begin{tabular}{p{0.29\linewidth}p{0.61\linewidth}}
\toprule
Dataset & Mathematical role\\
\midrule
Small binary actions & Literal actions through degree $16$, normal
axes, quotient maps, and the four local witness types.\\
Carrier data & Twelve noncritical actions, normal profiles, transition
data, and specified quotient-replacement maps.\\
Nonbinary and rank data & Finite action and normal-rank domains used
in the selected primitive and pair-top reductions.\\
Pair-top data & Concrete modules and constructive witnesses in the
count-twelve and count-sixteen domains.\\
Four- and six-block data & Constructive classifications and actual
quotient interfaces.\\
Zero-ternary frame data & Thirty complete binary module lattices and
thirty-eight literal odd-subgroup/normalizer witnesses.\\
Critical and audit data & Quadratic realization constants, critical
normalizers, and independent local chart controls.\\
\bottomrule
\end{tabular}
\end{center}
The data contracts specify literal point actions, normal subgroups,
source and target maps, and the original symmetric normalizers. A
quotient order or isomorphism-type name alone does not specify a
counting witness. For a proper subdirect replacement, the image group
and both its projection maps are part of the record.

For the critical quadratic constants of Lemma~\ref{lem:quadratic},
one can enumerate $\GL(2,2)$ and $\GL(4,2)$ and test the equation
$q(gv)=q(v)$ on every vector. This establishes the orders $2$ and $72$
for the displayed forms. The torsor argument in the lemma is what
extends these finite calculations to arbitrary source dimension.
Likewise the exact rational inequality \eqref{eq:positive-dual} is a
finite arithmetic certificate whose implication for all degree-$32$
sections is proved by the preceding incidence argument.

The classification libraries are pinned to
\textsf{TransGrp}~3.6.5 and \textsf{PrimGrp}~3.4.4
\cite{TransGrp365,PrimGrp344}; their mathematical construction and
scope are described in \cite{Hulpke2005,HoltRoyle2020,RoneyDougal2005Primitive}.
The recorded interpreter is GAP~4.13.1~\cite{GAP4131}.
Where a coverage proof uses an independent constructive generation
argument, the catalogue is a source of literal representatives, not
its sole completeness argument.

All reproduction commands are local. The file
\texttt{computations/README.md} documents the producers and independent
checkers, and \texttt{certificates/manifests/finite\_witnesses.json}
binds their finite input and output files. Hashes identify bytes; they
do not prove a mathematical theorem. The checks do not initialize or
invoke a proof assistant.

The same-source $C_3$ example, the inverse-complement reciprocal, the
proper-subdirect transport examples, and the binary duplicate factor
$1/4$ are regression conditions on the counting maps. Passing these
conditions alone does not establish an asymptotic theorem. The binary
coverage principles are proved in Lemma~\ref{lem:binary-coverage}; the
universal mixture and recurrence arguments are proved in
Theorems~\ref{thm:binary-mixture} and~\ref{thm:binary-recurrence}.
The complete nonbinary exhaustion is proved in
Theorem~\ref{thm:nonbinary-alphabet}; the ordinary counting inequality
and its two boundedness inductions are proved in Section~\ref{sec:assembly}.
The supplementary affine-nine and narrow socket data provide independent
local checks; the proof above uses the general frame and mixture arguments
for those domains.
